\section{A bound on uncountable base orders and categoricity on a tail}
\label{sec:categoricity}

Continue to work in ZFC with CH and with the class $K$ of
Section~\ref{sec:aec}. Recall the cardinals
\[
 \delta=\omega_1,\qquad \xi=2^\delta,
 \qquad \Lambda=\beth_{\xi^+}.
\]
Section~\ref{sec:endpoint} has supplied an object with uncountable $B$.
Its diagrams depend on the tuple and the coordinate, and their validity
has tuple-dependent countable exceptions. We first rule out a stronger
configuration: one fixed color for each arity passing every finite test.
A finite partition argument then forces precisely this configuration
if an uncountable-base object has at least $\Lambda$ index points.
This proves the universal bound $|I|<\Lambda$ for such objects.
Categoricity at and above $\Lambda$ will then follow by classifying the
remaining, countable-base objects through a matrix isomorphism calculation.

\subsection{The obstruction to a fixed pattern}

We first give the rank argument used to bound a diagram limit.
It is the ordinal-tagged form of the Mostowski collapse
\cite{Mostowski1949}; the construction and its range bound are proved
below. The following lemma concerns a set, with no countability
assumption on it or on its predecessor sets.

\begin{lemma}[Collapse of the tagged relation]\label{col:collapse}
Let \(Q\) be a set with an extensional binary relation \(R\) and
a function \(\rho:Q\to\delta\) such that
\[
                   yRx\ \Longrightarrow\ \rho(y)<\rho(x).
\]
Then there is an injection \(F:Q\to V_\delta\) satisfying
\[
                         F(x)=\{F(y):yRx\}.
\]
More precisely, \(F(x)\in V_{\rho(x)+1}\) for every \(x\in Q\).
\end{lemma}

\begin{proof}
Define \(F\) recursively in increasing order of the tag
\(\alpha<\delta\).  At tag \(\alpha\), all predecessors have
smaller tags, so their images have already been defined; the
displayed formula then defines the images of all elements
with tag \(\alpha\).  These are sets, since \(Q\) is a set.
Transfinite recursion and Replacement give a function on \(Q\).

We prove the stated bound by induction on \(\alpha=\rho(x)\).
If \(yRx\) and \(\beta=\rho(y)\), the induction hypothesis gives
\[
              F(y)\in V_{\beta+1}\subseteq V_\alpha,
              \qquad \beta<\alpha.
\]
Thus \(F(x)\subseteq V_\alpha\), and consequently
\(F(x)\in V_{\alpha+1}\).  At \(\alpha=0\) the predecessor
set is empty, so the same conclusion holds.  Since \(\delta\)
is a limit ordinal, \(\alpha+1<\delta\) for every
\(\alpha<\delta\), and the range of \(F\) lies in \(V_\delta\).

For injectivity, induct on \(\max(\rho(x),\rho(x'))\).
Suppose \(F(x)=F(x')\).  For each \(yRx\), equality of these
sets supplies \(y'Rx'\) with \(F(y)=F(y')\).  Both
\(\rho(y)<\rho(x)\) and \(\rho(y')<\rho(x')\), so the
induction hypothesis gives \(y=y'\).  This proves that every
predecessor of \(x\) is a predecessor of \(x'\), and the
reverse inclusion follows in the same way.  Extensionality
gives \(x=x'\).  At induction parameter zero both predecessor
sets are empty, so extensionality gives the conclusion directly.
\end{proof}

The proof uses the whole powerset \(V_{\alpha+1}=\mathcal P(V_\alpha)\).
Thus the bound remains valid when many local diagrams contribute
predecessors to the same element.

\begin{lemma}[No fixed pattern at all finite arities]\label{col:no-pattern}
There is no sequence \((s_j:j\ge1)\) of elements of \(S\) such that
\begin{equation}\label{col:fixed-pattern}
           T_k\bigl((s_{|J|})_{\varnothing\ne J\subseteq[k]}\bigr)
           \quad\text{holds for every }k\ge1.
\end{equation}
\end{lemma}

\begin{proof}
Suppose that such a sequence exists.  For each \(k\), let
\(\mathcal D_k\) be the diagram decoded from \(s_k\).
The tests assert that every \(\mathcal D_k\) is valid and that
\begin{equation}\label{col:pattern-restrictions}
 \mathcal D_k\restriction J=\mathcal D_{|J|}
       \qquad(\varnothing\ne J\subseteq[k]).
\end{equation}

Choose a linearly ordered set \(X\) with
\(|X|>|V_\delta|\).  Such an \(X\) can be the initial ordinal
of cardinality \(|V_\delta|^+\).  For each nonempty finite
\(A\subseteq X\), transport \(\mathcal D_{|A|}\) along the
increasing enumeration of \(A\), so that its labels become
\[
                 \{(U,i):\varnothing\ne U\subseteq A,\ i<\omega\}.
\]
Let \(Q_A\), \(R_A\), and \(\rho_A\) denote its quotient,
quotient relation, and tag function.
If \(A\subseteq B\) are nonempty and finite, the label inclusion
and Equation~\eqref{col:pattern-restrictions} give a map
\[
                         j_{A,B}:Q_A\longrightarrow Q_B.
\]
By Lemma~\ref{col:quotients}, these maps are injective, preserve
tags, preserve and reflect the relation, and satisfy
\[
                     j_{B,C}j_{A,B}=j_{A,C}
                     \qquad(A\subseteq B\subseteq C).
\]

We construct their limit explicitly.  Form the set
\[
       P=\{(A,a):\varnothing\ne A\subseteq X
                        \text{ is finite},\ a\in Q_A\}.
\]
This is a set-indexed union of countable sets.  Declare
\begin{equation}\label{col:limit-equivalence}
 (A,a)\sim(B,b)
 \quad\Longleftrightarrow\quad
 j_{A,A\cup B}(a)=j_{B,A\cup B}(b).
\end{equation}
Passing to a larger finite common stage preserves equality
and reflects it by injectivity.  The composition identities
therefore show that \(\sim\) is an equivalence relation:
in particular, two equalities witnessing transitivity can
both be read in the stage indexed by the union of their
three finite index sets.  Put \(Q=P/{\sim}\), and let
\(j_A:Q_A\to Q\) send \(a\) to the class of \((A,a)\).
Each \(j_A\) is injective, and
\[
                         j_Bj_{A,B}=j_A.
\]
Every finite collection of elements of \(Q\) lies in the
image of one \(j_A\), by taking a union of finitely many
index sets.

Define tags on the limit by
\(\rho(j_A(a))=\rho_A(a)\).  Define its relation using a
finite common stage: for \(a\in Q_A\), \(b\in Q_B\), and
\(C=A\cup B\), set
\begin{equation}\label{col:limit-relation}
 j_B(b)\,R\,j_A(a)
 \quad\Longleftrightarrow\quad
 j_{B,C}(b)\,R_C\,j_{A,C}(a).
\end{equation}
These definitions are independent of the representatives.
Indeed, representatives that describe the same limit points
become equal in a finite common stage by
Equation~\eqref{col:limit-equivalence}.  Taking the union of
that stage with all stages used in
Equation~\eqref{col:limit-relation} makes the relevant
representatives equal simultaneously.  Tag preservation
gives the same tag, and preservation and reflection of
\(R\) give the same truth value for the relation.
It also follows that each \(j_A\) preserves and reflects
\(R\), and that
\[
                         yRx\ \Longrightarrow\ \rho(y)<\rho(x)
                         \qquad(x,y\in Q).
\]

The limit relation is extensional.  To see this, take
distinct \(x,x'\in Q\) and put them in a common stage,
say \(x=j_A(a)\) and \(x'=j_A(a')\).  Injectivity implies
\(a\ne a'\).  Extensionality in \(Q_A\) gives a predecessor
distinguishing them.  After exchanging \(a,a'\) if needed,
there is \(b\in Q_A\) with
\[
                            bR_Aa,\qquad \neg(bR_Aa').
\]
Because \(j_A\) preserves and reflects the relation,
\(j_A(b)Rx\) and \(\neg(j_A(b)Rx')\).
This witness remains valid even if larger stages add
other predecessors.

For \(u\in X\), consider the element of \(Q_{\{u\}}\)
represented by the label \((\{u\},0)\), and let \(q_u\)
be its image in \(Q\).  If \(u\ne v\), condition
\textup{(iii)} of Definition~\ref{col:diagram}, applied
in the stage \(Q_{\{u,v\}}\), says that the two
corresponding elements are distinct.  Injectivity of
the stage map gives \(q_u\ne q_v\).  Thus
\[
                              |X|\le |Q|.
\]
Lemma~\ref{col:collapse} applies to the set \(Q\), with
its extensional relation and tags below \(\delta\).
It gives an injection into \(V_\delta\), whence
\[
                          |X|\le |Q|\le |V_\delta|.
\]
This contradicts the choice of \(X\).
\end{proof}

\begin{remark}[What the pattern quantifiers require]
\label{col:pattern-quantifiers}
Lemma~\ref{col:no-pattern} rules out one fixed sequence of
colors satisfying every finite test.  To obtain the
contradiction, it is enough to find, for each \(k\), a
separate ordered \(k\)-tuple whose \(j\)-element subsequences
all have color \(s_j\), and whose full subsequence test
passes.  Each such tuple certifies precisely the instance
of Equation~\eqref{col:fixed-pattern} for that \(k\).
The tuples need not be nested or come from a common
infinite homogeneous set: the fixed colors determine
the same decoded diagrams, and the tests determine
every increasing position restriction between them.
This is the form used below in Proposition~\ref{cat:bounded}.
\end{remark}

\subsection{The cardinal and finite partition estimates}

The finite-arity proof uses the end-homogeneous method of classical
partition calculus \cite{ErdosRado1956}. We give the particular estimate
and the subsequent selection of fixed colors explicitly, so their
cardinal bounds and quantifiers are part of the argument.

For a set $X$ and a positive integer $m$, write $[X]^m$ for its
$m$-element subsets. A subset $Y\subseteq X$ is \emph{homogeneous} for
a map $c:[X]^m\to C$ if $c$ is constant on $[Y]^m$.
Cardinals used as ordered sets below are identified with their initial
ordinals.

\begin{lemma}\label{cat:cardinals}
The cardinal $\Lambda$ is a strong limit and
\[
 \cf(\Lambda)=\xi^+>\xi.
\]
If $\zeta<\Lambda$ is infinite, then both $\xi^\zeta$ and
$(\xi^\zeta)^+$ are below $\Lambda$.
\end{lemma}

\begin{proof}
Put $\kappa=\xi^+$. This is a regular cardinal and a limit ordinal.
The beth sequence is strictly increasing and continuous at limit
ordinals, so
\[
 \Lambda=\sup_{\alpha<\kappa}\beth_\alpha.
\]
This gives $\cf(\Lambda)\le\kappa$. Conversely, a set of fewer than
$\kappa$ ordinals below $\Lambda$ can be bounded by beth values with
fewer than $\kappa$ indices. Regularity of $\kappa$ bounds these indices
below $\kappa$, and therefore bounds the original set below $\Lambda$.
This proves $\cf(\Lambda)=\kappa$.

Given an infinite cardinal $\nu<\Lambda$, choose $\alpha<\kappa$ with
$\nu\le\beth_\alpha$. Since $\alpha+1<\kappa$,
\[
 2^\nu\le 2^{\beth_\alpha}=\beth_{\alpha+1}<\Lambda.
\]
Thus $\Lambda$ is a strong limit. It is also a limit cardinal: for
every cardinal below it, a strictly larger beth value still lies below
it. In particular $\xi<\Lambda$, and for infinite $\zeta<\Lambda$,
\[
 \xi^\zeta\le (2^\xi)^\zeta
       =2^{\xi\cdot\zeta}<\Lambda.
\]
The last inequality uses $\xi\cdot\zeta<\Lambda$ and strong
limitness. The successor of $\xi^\zeta$ is also below the limit
cardinal $\Lambda$.
\end{proof}

\begin{lemma}[Finite partition estimate]\label{cat:finite-partition}
For every positive integer $m$ and every infinite cardinal
$\chi<\Lambda$, there is a cardinal $\mu<\Lambda$ such that every
map $c:[\mu]^m\to C$ with $|C|\le\xi$ has a homogeneous subset of
cardinality $\chi$.
\end{lemma}

\begin{proof}
By identifying the colors with elements of $\xi$, it suffices to use
the palette $\xi$. We induct on $m$, proving the statement for all
infinite targets $\chi<\Lambda$ at each step. When $m=1$, take
\[
 \mu=(\max(\xi,\chi))^+<\Lambda.
\]
If every color class had size less than $\chi$, their union would have
size at most $\xi\cdot\chi<\mu$. Thus some class has at least
$\chi$ elements, from which we select exactly $\chi$.

Suppose the statement is known for $m$, and fix a target
$\chi<\Lambda$. Choose an infinite cardinal $\zeta<\Lambda$ such
that every $\xi$-coloring of $[\zeta]^m$ has a homogeneous
$\chi$-element subset. Set
\[
 q=\xi^\zeta,\qquad \mu=q^+<\Lambda,
\]
where the last inequality follows from Lemma~\ref{cat:cardinals}.
Let $c:[\mu]^{m+1}\to\xi$. We seek $\zeta$ increasing representatives and a later point $x$ such
that the color of any $m+1$ representatives is unchanged when its last
point is replaced by $x$. Coloring an $m$-set by the color obtained
after adjoining $x$ then reduces to the induction hypothesis.
The residual-set tree below groups points by their profiles on the
named representatives along each branch. Its node count leaves an
unnamed $x$, whose residuals determine the required branch.

For each $\alpha<\zeta$, fix an injection
\begin{equation}\label{cat:profile-codes}
 \iota_\alpha:{}^{[\alpha+1]^m}\xi\longrightarrow q.
\end{equation}
Such an injection exists because $|[\alpha+1]^m|\le\zeta$.
Index potential tree nodes by sequences $s\in{}^\alpha q$, for
$\alpha\le\zeta$. Define their residual sets $X_s\subseteq\mu$
recursively as follows.
\begin{enumerate}[label=(\roman*)]
 \item At the root, put $X_\varnothing=\mu$.
 \item Suppose $s$ has length $\alpha<\zeta$ and $X_s$ is nonempty.
 Name its least point
 \[
  r_s=\min X_s.
 \]
 For $y\in X_s\setminus\{r_s\}$, define its profile
 $u_{s,y}\in{}^{[\alpha+1]^m}\xi$ by
 \begin{equation}\label{cat:profile}
  u_{s,y}(a)=c\bigl(\{r_{s\upharpoonright\beta}:\beta\in a\}
                    \cup\{y\}\bigr),
       \qquad a\in[\alpha+1]^m.
 \end{equation}
 For each $i<q$, put
 \[
  X_{s^\frown i}
   =\{y\in X_s\setminus\{r_s\}:\iota_\alpha(u_{s,y})=i\}.
 \]
 If $X_s$ is empty, set every immediate successor residual to the
 empty set.
 \item If $\gamma\le\zeta$ is a nonzero limit ordinal and
 $s\in{}^\gamma q$, put
 \begin{equation}\label{cat:limit-residual}
  X_s=\bigcap_{\alpha<\gamma}X_{s\upharpoonright\alpha}.
 \end{equation}
\end{enumerate}
Only nodes with nonempty residuals matter. Every such node has
nonempty ancestral residuals, so all representatives in
\eqref{cat:profile} are defined. Along a branch the representatives
increase strictly: each one is the least point of the current
residual, and successor residuals exclude that point. Consequently
the arguments of $c$ in \eqref{cat:profile} are $(m+1)$ distinct
points. Once $\alpha+1\ge m$, the index set $[\alpha+1]^m$ includes
subsets containing $\alpha$ itself, so the profile includes tests
using the just named point $r_s$.

The number of potential nodes, including those at limit levels, is
bounded by
\begin{equation}\label{cat:tree-count}
 \left|\bigcup_{\alpha\le\zeta}{}^\alpha q\right|
 \le |\zeta+1|\cdot q^\zeta
 =\zeta\cdot(\xi^\zeta)^\zeta
 =\xi^\zeta=q.
\end{equation}
Here $\zeta$ is infinite, $q\ge\zeta$, and
$\zeta\cdot\zeta=\zeta$. In particular at most $q$ points are
named at levels below $\zeta$. Since $\mu=q^+$, choose a point
$x\in\mu$ that is never named.

There is a unique branch of residuals containing $x$ through all
levels at most $\zeta$. Indeed, if $x\in X_s$ at a successor step,
then $x\ne r_s$, so exactly one of the profile classes contains it.
At a limit step, it belongs to the intersection
\eqref{cat:limit-residual} along the branch already constructed.
Thus the branch cannot disappear at a limit. Write its named
representatives as
\[
 r_\alpha\quad(\alpha<\zeta),
 \qquad P=\{r_\alpha:\alpha<\zeta\}.
\]
They increase strictly and are all less than $x$, so $|P|=\zeta$.

Define a coloring of the preceding arity by
\[
 d:[P]^m\longrightarrow\xi,
 \qquad d(A)=c(A\cup\{x\}).
\]
Consider indices $\alpha_0<\cdots<\alpha_m<\zeta$. At level
$\alpha_{m-1}$ of the branch, all the first $m$ representatives in
this list have appeared, including the just named one.
The last representative $r_{\alpha_m}$ and the point $x$ belong to
the same successor residual at that level. Their profiles therefore
agree on $\{\alpha_0,\ldots,\alpha_{m-1}\}$, giving
\begin{equation}\label{cat:arity-reduction}
 c(\{r_{\alpha_0},\ldots,r_{\alpha_m}\})
   =d(\{r_{\alpha_0},\ldots,r_{\alpha_{m-1}}\}).
\end{equation}
By the choice of $\zeta$, the map $d$ has a homogeneous subset
$H\subseteq P$ of cardinality $\chi$. Equation~\eqref{cat:arity-reduction}
shows that $H$ is homogeneous for $c$ as well. This proves the
induction step.
\end{proof}

\subsection{Extracting a forbidden pattern}

The next lemma combines the finite partition estimates without
asserting the existence of a common infinite homogeneous subset.

\begin{lemma}\label{cat:finite-patterns}
Let $X$ be a set with $|X|\ge\Lambda$. For each positive integer
$j$, let $c_j:[X]^j\to C_j$, where $|C_j|\le\xi$. There are colors
$p_j\in C_j$ such that for every positive integer $k$ there is
$A_k\in[X]^k$ satisfying
\[
 c_j(A)=p_j\qquad
 (1\le j\le k,\ A\in[A_k]^j).
\]
The sets $A_k$ need not be nested.
\end{lemma}

\begin{proof}
Choose the colors recursively. After choosing $p_1,\ldots,p_j$,
maintain the following invariant: for every infinite cardinal
$\chi<\Lambda$, there is a set $Y\subseteq X$ with
$\chi\le|Y|<\Lambda$ such that each $c_i$ for $1\le i\le j$
is constantly $p_i$ on $[Y]^i$. For $j=0$ this follows from
$|X|\ge\Lambda$.

Assume the invariant holds at $j$. For each infinite
$\chi<\Lambda$, Lemma~\ref{cat:finite-partition} supplies a source
cardinal $\mu_\chi<\Lambda$ for arity $j+1$ and target $\chi$.
Choose an existing sample of size at least $\mu_\chi$, restrict to
$\mu_\chi$ points, and apply the lemma to obtain a subset
$Z_\chi$ of size $\chi$ on which $c_{j+1}$ has a constant color
$q_\chi\in C_{j+1}$. All earlier colors are retained on this subset.

At least one color occurs as $q_\chi$ for target cardinals
$\chi$ unbounded below $\Lambda$. Otherwise, for every
$q\in C_{j+1}$ there would be a bound below $\Lambda$ for the
targets assigned color $q$. Since $|C_{j+1}|\le\xi<\cf(\Lambda)$,
the supremum of these bounds would still be below $\Lambda$.
This contradicts the fact that we made the construction for every
infinite target below the limit cardinal $\Lambda$. Choose such a
color as $p_{j+1}$; its samples establish the invariant at $j+1$.

After this recursion, fix a positive integer $k$. The invariant at
stage $k$ supplies an infinite sample homogeneous in the first $k$
chosen colors. Any $k$ points of that sample give $A_k$.
\end{proof}

\begin{proposition}\label{cat:bounded}
If $M\in K$ and $B^M$ is uncountable, then $|I^M|<\Lambda$.
\end{proposition}

\begin{proof}
Suppose instead that $|I^M|\ge\Lambda$. Fix a well-order of $I^M$,
used only in this proof, and fix row representatives $z_{t,n}$ for
every tuple $t$ and every $n<\omega$ in $M$. Let
\[
 c_t:B^M\longrightarrow S
\]
be the resulting color profile of $t$, as defined in
Section~\ref{sec:aec}. By CH, $|S|=\delta$, and
Lemma~\ref{aec:order-size} gives $|B^M|\le\delta$. Thus the palette
of whole profiles has size at most
\begin{equation}\label{cat:palette}
 |S^{B^M}|\le\delta^\delta=2^\delta=\xi.
\end{equation}
For the middle equality, $2^\delta\le\delta^\delta$ follows from
$2\le\delta$, while
$\delta^\delta\le(2^\delta)^\delta=2^\delta$ gives the reverse
inequality.

For each $j\ge1$, color a $j$-element subset of $I^M$ by the profile
of its increasing enumeration. Apply Lemma~\ref{cat:finite-patterns}
to these colorings, with every palette equal to $S^{B^M}$. Obtain
profiles $p_j\in S^{B^M}$ and, for each $k$, an increasing witness
tuple $t_k$ of length $k$ whose increasing $j$-subtuples all have
profile $p_j$, for $j\le k$.

For the fixed global row representatives, let $E_k\subseteq B^M$
be the countable exceptional set for the full object test of $t_k$.
By the definition of that test, it includes all required checks on
the nonempty subsequences of $t_k$. The set
\[
 E=\bigcup_{1\le k<\omega}E_k
\]
is countable. Choose $b\in B^M\setminus E$, possible because $B^M$
is uncountable, and put $s_j=p_j(b)$ for every $j\ge1$.
For each $k$, evaluation on $t_k$ proves
\[
 T_k\bigl((s_{|J|})_{\varnothing\ne J\subseteq\{1,\ldots,k\}}\bigr).
\]
This is the sequence forbidden by Lemma~\ref{col:no-pattern}.
The contradiction proves the bound.
\end{proof}

\begin{remark}\label{cat:separate-witnesses}
The witness tuples $t_k$ may be unrelated. At the chosen coordinate
$b$, the length-$k$ witness certifies validity of the one diagram
decoded from $s_k$, together with its restrictions along every
increasing inclusion of positions. Those fixed diagrams are exactly
the data to which Lemma~\ref{col:no-pattern} applies. The proof uses
only countably many full-tuple exceptional sets; it does not seek a
coordinate good for every tuple in $I^M$.
\end{remark}

\subsection{Isomorphisms when the base order is countable}

\begin{lemma}\label{cat:countable-base}
If $M,N\in K$ have countable base orders and $|I^M|=|I^N|$, then
$M\cong N$.
\end{lemma}

\begin{proof}
Fix a bijection $f:I^M\to I^N$. Both base orders are nonempty
countable dense linear orders without endpoints. There is an order
isomorphism $g:B^M\to B^N$, as follows. Enumerate the two orders and
build finite partial increasing bijections, alternately including the
next unused point of each enumeration in the domain and in the range.
At each step, density supplies a point between the finitely many
already matched neighbors, or absence of endpoints supplies a point
beyond them. The appropriate interval contains infinitely many
points, so the finitely many used points can be avoided. The union
is the required bijection $g$.

The map $g$ induces the group isomorphism
\[
 \widehat g:G(B^M)\longrightarrow G(B^N),
 \qquad u\longmapsto g[u],
\]
where elements of the coded groups are identified here with their
finite supports. Match the two copies of $G_*$ by their distinguished
elements. Enumerate $B^M$ as $(b_k:k<\omega)$; this enumeration need
not respect its order.

Fix one tuple $t$ of distinct points in $I^M$, and let $f(t)$ be its
coordinatewise image. Choose representatives
\[
 \begin{array}{ll}
 z_n\in Z^M_{t,n},&w_k\in W^M_{t,b_k},\\
 z'_n\in Z^N_{f(t),n},&w'_k\in W^N_{f(t),g(b_k)}.
 \end{array}
\]
Write their bit matrices as
\[
 a(n,k)=e^M(z_n,w_k),\qquad
 a'(n,k)=e^N(z'_n,w'_k),
\]
and let $d(n,k)=a(n,k)+a'(n,k)$, with addition in $\mathbb Z/2\mathbb Z$.
Define translators by
\begin{equation}\label{cat:triangle-supports}
 \begin{split}
 U_n&=\{g(b_k):k\le n\text{ and }d(n,k)=1\}\in G(B^N),\\
 V_k&=\{n:n<k\text{ and }d(n,k)=1\}\in G_*.
 \end{split}
\end{equation}
Every $U_n$ is finite, since it uses only the first $n+1$ columns;
every $V_k$ is finite, since it uses only the first $k$ rows. The two
regions in \eqref{cat:triangle-supports} partition the matrix, so
\begin{equation}\label{cat:triangle-equation}
 \boldsymbol 1_{g(b_k)\in U_n}
       +\boldsymbol 1_{n\in V_k}=d(n,k)\pmod2
       \qquad(n,k<\omega).
\end{equation}
Figure~\ref{fig:matrix-decomposition} illustrates this allocation of
finite supports.
\begin{figure}[tbp]
\centering
\begin{tikzpicture}[x=.52cm,y=.52cm,>=Latex,font=\small]
 \definecolor{rowshade}{RGB}{182,218,223}
 \definecolor{columnshade}{RGB}{249,221,167}
 \foreach \n in {0,...,6}{
  \foreach \k in {0,...,6}{
   \ifnum\k>\n
    \fill[columnshade] (\k,-\n) rectangle ++(1,-1);
   \else
    \fill[rowshade] (\k,-\n) rectangle ++(1,-1);
   \fi
  }
 }
 \draw[step=1,gray!60,thin] (0,0) grid (7,-7);
 \draw[->] (0,.9) -- (7.6,.9) node[right] {column \(k\)};
 \draw[->] (-1,0) -- (-1,-7.6) node[below] {row \(n\)};
 \foreach \k in {0,...,6}
  \node at (\k+.5,.35) {\(\k\)};
 \foreach \n in {0,...,6}
  \node at (-.4,-\n-.5) {\(\n\)};
 \draw[very thick,linkblue] (0,-4) rectangle (5,-5);
 \draw[very thick,orange!70!black] (5,0) rectangle (6,-5);
 \node[anchor=west,align=left] at (8,-1.2)
  {\(k\le n\): adjust in row \(n\)\\
   at most \(n+1\) positions};
 \node[anchor=west,align=left] at (8,-3.2)
  {\(n<k\): adjust in column \(k\)\\
   at most \(k\) positions};
 \node[anchor=west,align=left] at (8,-5.7)
  {Highlighted examples:\\
   row \(n=4\), column \(k=5\)};
 \node at (7.5,-6.5) {\(\cdots\)};
 \node at (6.5,-7.5) {\(\vdots\)};
\end{tikzpicture}
\caption{A finite portion of the countably infinite difference matrix.
The two regions partition its entries. Each row uses only the finitely
many positions \(k\le n\), and each column only the finitely many
positions \(n<k\); within these positions one shifts exactly where the
difference bit is one. Thus finite row and column translations remove
every difference matrix. The column enumeration need not respect the
dense order.}
\label{fig:matrix-decomposition}
\end{figure}

Set $\widetilde z'_n=z'_n+U_n$ and
$\widetilde w'_k=w'_k+V_k$. The action equation for the bit relation
and \eqref{cat:triangle-equation} give
\[
 e^N(\widetilde z'_n,\widetilde w'_k)
     =a'(n,k)+d(n,k)=a(n,k).
\]

Because every fiber is a torsor, its points have unique expressions
relative to its chosen representative. We therefore obtain
bijections of the row and column fibers by setting
\begin{equation}\label{cat:equivariant-maps}
 \begin{aligned}
 z_n+u&\longmapsto\widetilde z'_n+\widehat g(u)
                  &&(u\in G(B^M)),\\
 w_k+v&\longmapsto\widetilde w'_k+v
                  &&(v\in G_*).
 \end{aligned}
\end{equation}
These maps are equivariant and hence preserve the action graphs.
For arbitrary $u\in G(B^M)$ and $v\in G_*$, their target bit is
\[
 \begin{split}
 e^N(\widetilde z'_n+\widehat g(u),\widetilde w'_k+v)
 &=a(n,k)+\boldsymbol 1_{g(b_k)\in\widehat g(u)}
                +\boldsymbol 1_{n\in v}\\
 &=a(n,k)+\boldsymbol 1_{b_k\in u}+\boldsymbol 1_{n\in v}\\
 &=e^M(z_n+u,w_k+v).
 \end{split}
\]
Thus they preserve the entire bit relation, not just the bits at
representatives.

Carry out this construction separately for every tuple $t$. The row
and column fibers of distinct tuples are disjoint, and there are no
additional bit or action relations between different tuples. Together
with $f$, $g$, $\widehat g$, and the distinguished-element map on
$G_*$, these fiber maps give a bijection of the whole universes.
They preserve and reflect the sort predicates, tuple and base
projections, order, group incidence and operation graphs, action
graphs, distinguished elements, and bits. They therefore form an
isomorphism of the structures.
\end{proof}

\begin{theorem}\label{cat:tail}
The class $K$ is categorical in every cardinal $\mu\ge\Lambda$.
\end{theorem}

\begin{proof}
Fix such a cardinal $\mu$. By Lemma~\ref{aec:arrays}, choose a model
with $|I|=\mu$, base order $\Q$, the specified group torsors, and,
for example, zero bit sequences at the chosen representatives.
The countable base order makes the object test automatic.
Lemma~\ref{aec:size} gives cardinality exactly $\mu$, proving
existence.

Now let $M\in K$ have cardinality $\mu$. By
Lemmas~\ref{aec:order-size} and~\ref{aec:size},
\[
 |B^M|\le\delta<\Lambda\le\mu,
 \qquad
 \mu=\max(|I^M|,|B^M|,\aleph_0).
\]
It follows that $|I^M|=\mu$. Proposition~\ref{cat:bounded} therefore
forces $B^M$ to be countable. Any two models of size $\mu$ now have
equally sized index sorts and countable base orders, so
Lemma~\ref{cat:countable-base} makes them isomorphic. This proves
existence of exactly one isomorphism type at every $\mu\ge\Lambda$,
including $\mu=\Lambda$ itself.
\end{proof}
